% Part: computability
% Chapter: computability-theory

\documentclass[../../../include/open-logic-chapter]{subfiles}

\begin{document}

\olchapter{cmp}{thy}{గణనీయతా సిద్ధాంతం}

\begin{editorial}
  ఈ అధ్యాయంలోని విషయాన్ని సమీక్షించి విస్తరించాలి. ముఖ్యంగా, ఇంకా
  సాధనలు లేవు.
\end{editorial}

\olimport{introduction}

\olimport{coding-computations}

\olimport{normal-form}

\olimport{s-m-n}

\olimport{universal-part-function}

\olimport{no-universal-function}

\olimport{halting-problem}

\olimport{russells-paradox}

\olimport{computable-sets}

\olimport{ce-sets}

\olimport{equiv-ce-defs}

\olimport{non-comp-set}

\olimport{ce-closed-cup-cap}

\olimport{complement-ce}

\olimport{reducibility}

\olimport{prop-reduce}

\olimport{complete-ce-sets}

\olimport{k-1}

\olimport{total}

\olimport{rice-theorem}

\olimport{fixed-point-thm}

\olimport{application-fixed-point}

\olimport{def-functions-self-reference}

\OLEndChapterHook

\end{document}
